%% kapa: exam=1 round=36 year=2025 period=1 form=A subject=real_estate subject_ko=부동산학원론 seq=9 number=89
%% source: Q-Net 감정평가사 기출문제 (cst00309) artlSeq=5231525 file=raw/1008/5231525/2025년 제36회 감정평가사 1차시험 1교시 문제지 원본.pdf
%% conversion: kordoc 4.24.51 -> tools/gen1.py (auto); review=proofread-sonnet; flags=none
\begin{question}{89}
지대이론에 관한 설명으로 옳은 것은 모두 몇 개인가?
\begin{framedtext}
○ 튀넨(J.H. von Thünen)은 자연조건이 동일한 고립국을 가정하여 상업활동의 공간적 분포를 통한 토지이용을 설명한다. \par
○ 리카도(D. Ricardo)는 각 토지마다 다른 비옥도의 차이와 생산요소 투입에 따라 한계생산성이 증가하는 수확체감현상을 적용한다. \par
○ 마샬(A. Marshall)은 생산요소에 귀속되는 소득으로서 생산품의 총판매수익에서 가변비용을 제외한 잉여분을 절대지대라고 주장한다. \par
○ 알론소(W. Alonso)는 해당 토지의 지대를 지대입찰과정에서 토지이용자가 지불하고자 하는 최고 지불용의액으로서 초과이윤이 0(zero)이 되는 지대로 보았다. \par
○ 해리스(C. Harris)와 울만(E. Ullman)은 토지이용자가 공간의 마찰비용으로 지대와 교통비를 함께 지불한다고 보았다.
\end{framedtext}
\begin{choices}
  \item 1개
  \item 2개
  \item 3개
  \item 4개
  \item 5개
\end{choices}
\end{question}
\answer{1}
